\documentclass[11pt,a4paper]{article}

\usepackage[T1]{fontenc}
\usepackage[utf8]{inputenc}
\usepackage{lmodern}
\usepackage[margin=2.6cm]{geometry}
\usepackage{amsmath,amssymb}
\usepackage{booktabs}
\usepackage{tabularx}
\usepackage{enumitem}
\usepackage{xcolor}
\usepackage[most]{tcolorbox}
\usepackage[hidelinks]{hyperref}

\definecolor{navy}{RGB}{0,0,110}
\definecolor{teamblue}{RGB}{0,70,140}
\definecolor{indivred}{RGB}{190,0,90}
\definecolor{boxgrey}{RGB}{140,140,140}

\setlength{\parindent}{1.5em}
\setlength{\parskip}{0.3em}
\setlength{\emergencystretch}{2em}

% Question symbols
\newcommand{\team}{\textcolor{teamblue}{$\bullet$}}
\newcommand{\indiv}{\textcolor{indivred}{$\blacklozenge$}}
\newcommand{\opt}{$\circ$}

\newcommand{\E}{\mathbb{E}^{\mathbb{Q}}}
\newcommand{\ind}[1]{\mathbf{1}_{\{#1\}}}

% Boxes
\newtcolorbox{navybox}[1]{enhanced, breakable, colback=navy!4, colframe=navy,
  fonttitle=\bfseries, title={#1}, arc=1mm, boxrule=0.6pt}
\newtcolorbox{greybox}[1]{enhanced, breakable, colback=black!4, colframe=boxgrey,
  fonttitle=\bfseries, title={#1}, arc=1mm, boxrule=0.6pt}

% Question list
\newlist{questions}{enumerate}{1}
\setlist[questions]{label=\textbf{Q\arabic*.}, leftmargin=3em, itemsep=0.6em}
\newlist{dquestions}{enumerate}{1}
\setlist[dquestions]{label=\textbf{D\arabic*.}, leftmargin=3em, itemsep=0.6em}

\begin{document}

\begin{center}
  {\LARGE\bfseries Advanced Model Risk Management}\\[0.4em]
  {\large Tutorials 1--2 --- The Phoenix Autocall as a Case Study}\\[2.5em]
  Two 90-minute sessions --- Notebook: \texttt{TD\_model\_risk.ipynb} (posted on the ENT)\\[0.6em]
  \rule{0.7\textwidth}{0.4pt}
\end{center}

\section*{Foreword: why this product}

Model risk is hard to teach because most examples are degenerate: either the model is obviously
wrong, or the product is so simple that every model agrees. The Phoenix autocall avoids both traps.
It is one of the few instruments for which \emph{several models calibrate to exactly the same market data and
still give different prices.} The divergence therefore comes neither from an implementation error nor
from a bad set of parameters: it is structural, and that is precisely the definition of model risk.

The tutorial follows seven steps. We start from the \textbf{payoff} and its decomposition, infer \textbf{which
risks} the product carries, then \textbf{which model} is acceptable, and then turn the question around by
judging models no longer on their prices but on \textbf{how faithful their sensitivities are}. We end with
two warnings: faithful greeks are not a sufficient criterion, and enriching a model is not always the
right answer.

\section*{How the tutorials work}

\noindent\textbf{Vocabulary.} TD = tutorial session: TD 1 and TD 2 are the two 90-minute sessions.\\
\textbf{Teams.} You work in \textbf{teams of 4 to 5}. Each team has its own \textbf{term sheet}: a coupon barrier $B_c$
and a protection barrier $B_p$, given by your group and team number in the team table (Excel answer
sheet, tab \emph{Teams}, and notebook \S0). Every number you report is for your own term sheet.\\
\textbf{The parts are independent.} Sections 1 to 8 and Block D can be read and answered in any order:
each question recalls the result it needs, and the numbers it relies on come from the notebook. The
only prerequisite is to run notebook \S0 with your group and team number.\\
\textbf{Who answers what.} Each question carries a symbol:

\begin{center}
\small
\begin{tabular}{@{}c l p{6.2cm}@{}}
\toprule
\team  & \textbf{Team}                  & Answered by the team in the team answer sheet (Excel), including its \emph{Decision} tab. \\
\indiv & \textbf{Individual, mandatory} & Answered by \emph{each student alone}, in the individual answer sheet (Excel). \\
\opt   & \textbf{Optional}              & For those who want to go further; not graded. \\
\bottomrule
\end{tabular}
\end{center}

\noindent\textbf{Before TD 1: read the subject.} Please take the time to read this subject before the first session:
at least Section 1 (the product, the definitions, the parameters and the reminders (R1)--(R9)), and
skim the questions. Also run the notebook once at home ($\sim$3 minutes) to check that it works. You
will lose much less time in class.

\begin{center}
\small
\begin{tabularx}{\textwidth}{@{}l >{\raggedright\arraybackslash}X >{\raggedright\arraybackslash}X@{}}
\toprule
\textbf{Phase} & \textbf{Questions} & \textbf{What happens} \\
\midrule
\textbf{Before} & Read the subject (at least Section 1) + run the notebook once & Not graded. Check at home that the notebook runs. \\[0.4em]
\textbf{TD 1} (90 min) & Q10, 11, 12, 13, 16, 14 & The model gap, on your term sheet. Fill in part 1 of the answer sheet (kept by the team). \\[0.4em]
\textbf{TD 2} (90 min) & Q18, 22--26 + Block D (D1--D3) + \emph{Decision} tab & What to do with the gap; a pricer written by an AI; your validation decision. Teams that have finished may start the individual part on their own. \\[0.4em]
\textbf{After} & Team: finish the answer sheet (one week). Individual: all \indiv\ questions (end of October) & Team answer sheet and notebook; individual answer sheet (+ optional PDF appendix). \\
\bottomrule
\end{tabularx}
\end{center}

\begin{navybox}{Roadmap of the two sessions (minutes)}
\small
\textbf{TD 1 (whole team). 0--5}: set \texttt{GROUP} and \texttt{TEAM} in \S0, run the \emph{whole} notebook ($\sim$3 min), and meanwhile
read Section 1.3 and, in Section 1.4, (R1), (R4) and (R6); designate your \textbf{Block D lead} and enter his
or her email in the answer sheet (tab \emph{Team}). \textbf{5--15}: Q10. \textbf{15--28}: Q11. \textbf{28--52}: Q12--13. \textbf{52--72}: Q16.
\textbf{72--82}: Q14 (shorten it if you are running late). \textbf{82--90}: answer sheet, part 1 (Excel template). Nothing is
submitted at the end of TD 1.

\textbf{TD 2 (the team splits in two). 0--3}: recap of model M4 by the instructor. \emph{Block D sub-team (the lead
and one or two members).} \textbf{3--13}: test plan (D1), handed in \emph{before} receiving the module; the instructor then
emails your team's module to the Block D lead, who shares it with the team. \textbf{13--50}: run the tests (D2)
and measure the impact of each error (D3). \emph{Rest of the team.} \textbf{3--15}: Q18. \textbf{15--30}: Q22, Q23, Q24. \textbf{30--50}:
Q25--26 (notebook \S7, run by you). \emph{Whole team.} \textbf{50--75}: the \emph{Decision} tab (Section 10). \textbf{75--90}: complete
the answer sheet; teams that have finished start the individual part on their own. Nothing is submitted at
the end of a session: see the deadlines below.

\textbf{Precision.} Every number you report comes with the number of paths and a standard error (printed by
the notebook for prices and break-even coupons), or failing that a second run with another seed.\\
\textbf{Setup.} Python $\geq$ 3.9 with numpy, scipy and matplotlib (recent numpy: see the first cell of the notebook).
\end{navybox}

\noindent\textbf{Grading: 50\,\% team, 50\,\% individual.} There is no oral presentation. \emph{Team work} (50\,\%): the
team answer sheet \texttt{TD\_answer\_sheet.xlsx} --- numbers, precision and Block D tests (25\,\%); \emph{Decision} tab (25\,\%). \emph{Individual work} (50\,\%): answers to the \indiv\ questions in the individual answer
sheet \texttt{TD\_individual\_answers.xlsx}, written alone; long derivations may go in an optional PDF appendix.
Everything is submitted on \textbf{Moodle} (or by email to the instructor if Moodle is unavailable). Groups
IF3 and IF2, whose two tutorials run back to back on the same evening, are graded with the same
leniency for this atypical schedule. The detailed grid is in the \textbf{evaluation sheet}.

\begin{center}
\small
\begin{tabular}{@{}l l l@{}}
\toprule
\textbf{Group} & \textbf{Team answer sheet + notebook} & \textbf{Individual answer sheet (+ PDF appendix)} \\
               & (one week after TD 2)                  & (end of the month) \\
\midrule
IF3 & Fri, Oct 9, 11:59 pm  & Sat, Oct 31, 11:59 pm \\
IF2 & Thu, Oct 22, 11:59 pm & Sat, Oct 31, 11:59 pm \\
IF1 & Fri, Oct 23, 11:59 pm & Sat, Oct 31, 11:59 pm \\
\bottomrule
\end{tabular}
\end{center}


Several questions explicitly recall the result to use: they should not block anyone. The technical terms
of Section 1.2 (EMTN, \emph{worst-of}, \emph{funding}\ldots) are not needed in class.

%======================================================================
\section{The product and what is at stake}

\subsection{Mechanics}

A \emph{Phoenix autocall} is a debt security (typically an EMTN) whose cash flows are linked to an equity
underlying and observed on fixed dates. Three building blocks, checked on each date $t_i$:
\begin{itemize}
  \item \textbf{Conditional coupon}: paid if $S_{t_i} \geq B_c S_0$ (coupon barrier, e.g.\ 70\,\%), with or without a \emph{memory}
  feature (catch-up of missed coupons).
  \item \textbf{Early redemption} (\emph{autocall}): if $S_{t_i} \geq B_a S_0$ (e.g.\ 100\,\%), the note is redeemed at par and all
  future cash flows disappear.
  \item \textbf{Conditional protection}: at maturity, redemption at par unless $S_T < B_p S_0$ (e.g.\ 60\,\%), in which
  case the investor bears the underlying's performance.
\end{itemize}

\begin{greybox}{Definitions used throughout}
\small
\textbf{At par.} At 100\,\% of the notional. The note is issued at par (the investor pays 100) and, when it is called
or protected, redeemed at par (the investor receives 100). \emph{Below par} or \emph{above par} means a price under or
over 100\,\%.\\
\textbf{Expected life (expected duration).} The redemption date $\tau$ of an autocall is random: it depends on the
path of the underlying. Its \emph{expected life} is $\E[\tau]$, the average time the note stays alive. A plain bond has a
known life; an autocall's life shortens when early redemption becomes more likely.\\
\textbf{Separable instrument.} A feature of a contract is \emph{separable} when it could be detached and traded on
its own, as a stand-alone contract with its own cash flows, without changing the rest of the contract (the
accounting notion of an embedded derivative that can be split from its host).\\
\textbf{Structuring equation.} The identity that, at issue, splits the issue price (par) among the building blocks
of the product and the fees kept by the bank: it shows what finances the coupon.\\
\textbf{Observed at each close.} A barrier checked against the official \emph{closing price} of the underlying on every
trading day, rather than only at maturity (European barrier) or at every instant during the day (continuous
barrier).
\end{greybox}

\subsection{Economics}

\begin{greybox}{What a Phoenix really is}
The coupon is not a yield: it is the price of an option sold by the investor. The product is a
\textbf{bet that the market will not fall}, never a bet that it will rise --- there is no participation in
the upside. The distribution of outcomes is highly skewed: many small gains, few losses but very
heavy ones. The median is flattering; the mean much less so.
\end{greybox}

Five points shape the economics of the product and will come back in the tutorial.
\begin{enumerate}
  \item \textbf{Where the budget comes from.} The coupon is financed by the premium of the put sold by the
  investor and by the issuer's \emph{funding} benefit (the zero-coupon bond is issued on the issuer's credit
  curve). Rate levels and credit spreads therefore drive how generous the coupon can be as much as
  the level of volatility does.
  \item \textbf{Early redemption shortens the expected life.} The most likely scenario (a stable or rising
  market) is also the one that pays least in absolute terms: an early call, one or two coupons,
  reinvestment risk.
  \item \textbf{The choice of underlying dominates everything else.} A \emph{worst-of} on several uncorrelated
  stocks, or a \emph{decrement} index (a fixed synthetic dividend subtracted, hence a lower forward), raise
  the displayed coupon far more than any barrier setting. They raise the risk in the same proportion.
  \item \textbf{The issuer's credit risk is real.} ``Capital protection'' is only as good as the issuer's signature.
  \item \textbf{The secondary market is a monopoly.} The issuer is usually the only market maker, and
  distribution fees, paid upfront, are amortized in the exit price.
\end{enumerate}

\subsection{Parameters used in the tutorial}

\begin{center}
\begin{tabular}{@{}l l@{}}
\toprule
Underlying, initial spot $S_0$ & single stock, 100 \\
Maturity $T$ & 3 years \\
Observation dates & semi-annual: 0.5, 1.0, \ldots, 3.0 \\
Autocall barrier $B_a$ & 100\,\% \\
Coupon barrier $B_c$ & 70\,\% (reference case; your team: 65, 70 or 75\,\%) \\
Protection barrier $B_p$ & 60\,\% (reference case; your team: 50, 55 or 65\,\%) \\
Memory feature & no (unless stated otherwise) \\
Risk-free rate $r$, dividend yield $q$ & 2\,\%, 2\,\% \\
Issue price & at par (100\,\%) \\
\bottomrule
\end{tabular}
\end{center}

\noindent\textbf{Your term sheet.} Wherever the subject says 70\,\% (coupon) or 60\,\% (protection), read the barriers
of your team's term sheet; in the notebook, set \texttt{GROUP} and \texttt{TEAM} in the parameter cell of \S0 and the
barriers follow. The numbers quoted in this subject are those of the reference case.

The \textbf{market} is represented by a reference Heston model ($v_0 = \theta = 4\,\%$, $\kappa = 1.5$, $\xi = 0.6$, $\rho = -0.7$),
from which the whole implied volatility surface is derived. The 3-year ATM implied volatility is about
17\,\%, with a pronounced skew.

\subsection{Reminders, notation and results taken as given}

This section gathers everything used in the tutorial. \textbf{None of these results needs to be proved
again}: they are provided to be used.

\paragraph{(R1) Risk-neutral price and stopping time.} Under a risk-neutral measure $\mathbb{Q}$, writing $R_i =
S_{t_i}/S_0$ and
\[
  \tau \;=\; \inf\{\, t_i,\; i<n \;:\; R_i \geq B_a \,\}, \qquad \inf\emptyset := t_n,
\]
the price of the Phoenix reads
\[
  V_0 \;=\; \E\!\left[\, \sum_{i=1}^{n} e^{-r t_i}\, c\, \ind{R_i \geq B_c}\, \ind{t_i \leq \tau} \;+\; e^{-r\tau}\, \Pi_\tau \right],
  \quad
  \Pi_\tau =
  \begin{cases}
    1 & \tau < t_n \\
    1 - (1 - R_n)\,\ind{R_n < B_p} & \tau = t_n.
  \end{cases}
\]

\paragraph{(R2) The four models.} All under $\mathbb{Q}$, with $S_0$ and $r, q$ given.
\begin{gather*}
  \textbf{M1}:\;\; \frac{dS_t}{S_t} = (r-q)\,dt + \sigma_{\mathrm{ATM}}\, dW_t
  \qquad\qquad
  \textbf{M2}:\;\; \frac{dS_t}{S_t} = (r-q)\,dt + \sigma_{\mathrm{loc}}(S_t,t)\, dW_t \\[0.6em]
  \textbf{M3}:\;\; \frac{dS_t}{S_t} = (r-q)\,dt + \sqrt{v_t}\, dW^1_t, \quad
  dv_t = \kappa(\theta - v_t)\,dt + \xi\sqrt{v_t}\, dW^2_t, \quad
  d\langle W^1, W^2\rangle_t = \rho\, dt. \\[0.6em]
  \textbf{M4 (local-stochastic)}:\;\; \frac{dS_t}{S_t} = (r-q)\,dt + L(S_t,t)\sqrt{v_t}\, dW^1_t, \\
  dv_t = \kappa(\theta - v_t)\,dt + \boldsymbol{\eta}\,\xi\sqrt{v_t}\, dW^2_t,
\end{gather*}
where $L$ is the \textbf{leverage function} and $\eta \in [0,1]$ the \textbf{mixing fraction}. M4 is the market standard for
autocalls; it was only mentioned in the lectures. \emph{What to remember:} at each date, $L$ is chosen so that
$L^2(S,t)\,\E[v_t \mid S_t = S] = \sigma^2_{\mathrm{loc}}(S,t)$; by Gyöngy (R4), every value of $\eta$ therefore reproduces \emph{all} vanillas.
$\eta$ sets the share of stochastic volatility: $\eta = 0$ gives back local volatility M2, $\eta = 1$ stochastic volatility
M3. Vanillas do not see $\eta$; the forward skew does. The notebook provides M4 already calibrated.

\paragraph{(R3) Dupire's formula.} In total implied variance $w(y,T) = \sigma^2_{\mathrm{imp}}(y,T)\,T$ with $y = \ln(K/F_T)$:
\[
  \sigma^2_{\mathrm{loc}}(y,T) \;=\;
  \frac{\partial_T w}
  {1 - \frac{y}{w}\,\partial_y w + \frac{1}{4}\left(-\frac{1}{4} - \frac{1}{w} + \frac{y^2}{w^2}\right)(\partial_y w)^2 + \frac{1}{2}\,\partial_{yy} w}.
\]

\paragraph{(R4) Gyöngy's theorem (1986), taken as given.} Let $S$ be a stochastic volatility diffusion.
Setting $\sigma^2_{\mathrm{loc}}(K,T) = \E[v_T \mid S_T = K]$, the process $\widetilde S$ solving $d\widetilde S_t/\widetilde S_t = (r-q)\,dt + \sigma_{\mathrm{loc}}(\widetilde S_t, t)\,dW_t$ has,
\emph{for each} $T$, the same marginal law as $S_T$. Nothing is claimed about the joint laws of $(S_{T_1},\ldots,S_{T_k})$.

\paragraph{(R5) Completeness criterion, taken as given.} A market driven by $d$ independent Brownian
motions is complete if $d$ traded assets (besides the numeraire) are available whose volatility matrix is
invertible. In this tutorial, unless stated otherwise, we hedge with the \emph{underlying alone} (and cash).
Otherwise, the martingale measure is not unique and some payoffs cannot be replicated.

\paragraph{(R6) Forward variance and forward smile.} The forward variance is $\E\big[\int_{T_1}^{T_2} v_s\,ds \,\big|\, \mathcal{F}_0\big]$. The
\textbf{forward smile} $T_1 \to T_2$ is the implied volatility surface of \emph{forward-starting} options, with payoff
$\big(S_{T_2}/S_{T_1} - k\big)^+$. Its slope is the \textbf{forward skew}. It depends on the law of $S_{T_2}$ \emph{given} $S_{T_1}$ (the \emph{transition
kernel}), which calibration to vanillas does not constrain.

\paragraph{(R7) The three valuation layers, and what they correct.}

\begin{center}
\small
\begin{tabularx}{\textwidth}{@{}>{\raggedright\arraybackslash}X >{\raggedright\arraybackslash}X >{\raggedright\arraybackslash}X >{\raggedright\arraybackslash}X@{}}
\toprule
\textbf{Layer} & \textbf{Framework} & \textbf{What it corrects} & \textbf{Where it is booked} \\
\midrule
Front-office valuation & internal (management) & nothing: it is the desk model's \emph{mid} price & management P\&L, greeks \\[0.3em]
Fair value adjustments & \textbf{accounting}: IFRS 13, IFRS 9 & the gap between the \emph{mid} price and the \emph{exit price} of an actual position & income statement and balance sheet (fair value) \\[0.3em]
Prudent valuation (AVA) & \textbf{prudential}: CRR art.\ 105, EBA RTS & the gap between fair value and a \emph{prudent} value (90\,\% confidence) & deduction from CET1 capital \\
\bottomrule
\end{tabularx}
\end{center}

\textbf{IPV is not a reserve.} \emph{Independent Price Verification} (CRR art.\ 105(8)) is a \textbf{control}: a periodic
check --- at least monthly --- of the accuracy of the front office's parameters and marks against
independent sources. It produces \emph{differences}, which may trigger an adjustment; it is not one.

\textbf{Usual fair value adjustments (accounting):} bid-offer / close-out cost; parameter uncertainty
(unobservable inputs); model uncertainty; concentration and liquidity; XVA (CVA, DVA, FVA); and,
for Level 3 positions, \textbf{deferral of day-one P\&L} (IFRS 9) --- which is not a reserve but an accounting
rule.

\textbf{AVA categories (prudential):} market price uncertainty --- close-out costs --- \textbf{model risk} ---
concentrated positions --- future administrative costs --- early termination --- operational risk ---
investing and funding costs --- unearned credit spreads. AVAs \textbf{come on top of} the accounting
adjustments already booked: they are added to them, and a risk already covered by an accounting
adjustment is not counted a second time.

\textbf{Three kinds of uncertainty not to be confused:} \emph{(i)} a poorly observed \textbf{market} parameter ---
the quantity exists and trades, but without a reliable consensus (long-dated dividends, repo, implied
correlation of a basket); \emph{(ii)} a \textbf{model} parameter --- the quantity exists only inside the model and no
data can reveal it ($\eta$); \emph{(iii)} \textbf{specification} risk --- the choice of the model family itself. The three call
for different treatments and governance.

\paragraph{(R8) Vocabulary.} \emph{Vega convexity} (or \emph{volga}): $\partial^2 V/\partial\sigma^2$, zero for an at-the-money vanilla, nonzero
elsewhere. --- \emph{Vega buckets}: breakdown of vega by maturity bucket. --- \emph{Day-one P\&L}: profit recognized
at inception, equal to the gap between the transaction price and the model value. --- \emph{AVA} (\emph{Additional
Valuation Adjustment}): see (R7). --- \emph{Cross gamma}: second-order cross sensitivity $\partial^2 V/\partial x\,\partial y$ between
two risk factors. --- \emph{Gap risk}: loss caused by a move of the underlying across a level where the payoff is
discontinuous (a barrier), between two rebalancings or through a jump. Near the barrier and maturity,
the delta becomes so large that the hedge cannot follow: the loss is of the order of the payoff's jump,
whatever the model. --- \emph{Delta} (notebook, \S7): change in value for a 100\,\% rise of the spot, as a fraction
of the notional, i.e.\ $S\,\partial V/\partial S$; a delta of 20 means that a 1\,\% move of the spot changes the value by
20\,\% of the notional.

\paragraph{(R9) What a model is (supervisory definition).} A \emph{model} is a quantitative method, system
or approach that applies statistical, economic, financial or mathematical theories, techniques and
assumptions to process input data into quantitative estimates. It has three components: an input
component, a processing component that turns inputs into estimates, and a reporting component.
Who or what wrote the code plays no role in this definition. Model changes are graded by materiality,
as in common bank policies (Session 5): \textbf{I --- minor} (routine recalibration, data refresh, cosmetic code
change): notification and automated tests; \textbf{II --- material} (new product in scope, new calibration
set, library upgrade): targeted review before use; \textbf{III --- major} (new methodology, new model, new
numerical scheme): full validation before use.

\begin{navybox}{Convention: the break-even coupon}
Coupons are paid \textbf{semi-annually}. The \textbf{break-even coupon} is the semi-annual coupon $c$ such
that the value of the product equals par. Since the value is affine in $c$,
\[
  c \;=\; \frac{1 - \text{(redemption leg)}}{\text{annuity}}, \qquad
  \text{annuity} = \sum_i \E\big[e^{-r t_i}\, \mathbf{1}_{\{\text{coupon paid at } t_i\}}\big].
\]
It is quoted \textbf{per year}: the annual coupon is $2c$, as on a term sheet (``6\,\% p.a., paid semi-annually'').
Coupons are paid out, not reinvested, so there is no compounding: not $(1+c)^2 - 1$. Every model
gap is expressed in \textbf{basis points of annual coupon} (twice the semi-annual gap), the one unit
that speaks to the distributor and to the trader alike.
\end{navybox}

%======================================================================
\section{Decomposing the payoff}

\begin{questions}
  \item \indiv\ Write the cash flow of the product on date $t_i$ as a function of $S_{t_i}$ and of the survival indicator
  $\mathbf{1}_{\{\text{not called before } t_i\}}$. Deduce a decomposition of the Phoenix into four elementary building blocks:
  a zero-coupon bond, a sold barrier put, a strip of bought digitals, and a \emph{knock-out} condition.

  \item \indiv\ \textbf{Show that} the price of the product reads as in (R1), explaining the role of each indicator.
  \textbf{Deduce that} the payoff is not a function of $S_T$ alone: state the variables it depends on. \emph{(The
  result to establish is given in (R1); the task is to justify it term by term.)}

  \item \indiv\ Explain why the autocall clause is not a separable instrument (see the definitions in Section 1.1)
  but a condition on \emph{all} future cash flows. What does this imply for the expected life of the
  product?

  \item \indiv\ Write the structuring equation linking the coupon, the put premium, the cost of the digital
  strip, the funding benefit, the structuring margin and the distribution fee.

  \item \indiv\ The put sold here has a \textbf{European barrier}: the protection barrier is only observed at maturity.
  Write its payoff. How does it differ from a \emph{down-and-in} put with an American barrier
  (observed continuously, or at each close)? Which of the two is more expensive, and why did the
  European market switch massively to the first one after 2008?
\end{questions}

%======================================================================
\section{A taxonomy of model risk and a decision ladder}

We use the usual three-way taxonomy: \textbf{(i)} the model is badly implemented; \textbf{(ii)} the model is correct
but fed with bad inputs; \textbf{(iii)} the chosen model is not the right one for this derivative. Form (i) is the
subject of Block D (Section 9); everywhere else, implementations are assumed to be sound.

\begin{questions}[resume]
  \item \indiv\ For each of the three forms, give a concrete example drawn from the Phoenix.

  \item \indiv\ Consider the following \textbf{decision ladder}, which ranks payoffs by increasing demands on the
  model:
  \begin{center}
  \small
  \begin{tabular}{@{}c l l@{}}
  \toprule
  Level & Nature of the payoff & Minimal model \\
  \midrule
  A & decomposable into vanillas & Black--Scholes, one implied vol per strike \\
  B & sensitive to the skew, not decomposable & properly calibrated local vol \\
  C & vega convexity, forward skew & stochastic volatility \\
  \bottomrule
  \end{tabular}
  \end{center}
  Rank: (a) a call spread; (b) an Asian option; (c) a European digital; (d) a European-barrier
  put; (e) a down-and-in put with a continuous barrier; (f) a cliquet with local floors and caps;
  (g) the full Phoenix. Justify each answer in one sentence.

  \item \indiv\ Why can a call spread be priced in Black--Scholes although it is sensitive to the skew? What
  in the Phoenix prevents the same reasoning?

  \item \indiv\ \textbf{Using} the dynamics (R2) and the completeness criterion (R5): for each of the three models,
  give the number of Brownian motions and the number of traded assets when hedging with
  the \emph{underlying alone}. \textbf{Deduce that} M1 and M2 then define complete markets, and M3 an
  incomplete one. What does this imply in practice for what one can expect from a delta hedge
  under M3? What would adding vanillas to the hedging instruments change?
\end{questions}

%======================================================================
\section{Calibration \emph{is} a model choice}

This is the heart of the tutorial. We have three models: \textbf{(M1)} Black--Scholes with a flat volatility set
at the 3-year ATM; \textbf{(M2)} Dupire local volatility calibrated to the whole surface; \textbf{(M3)} the reference
stochastic volatility model. By construction, M2 and M3 reproduce \emph{exactly the same} vanilla prices.

\begin{questions}[resume]
  \item \team\ \textbf{Using Gyöngy's theorem (R4)}, show that M2 and M3 give the \emph{same} value to any payoff
  of the form $g(S_T)$. Check this result numerically \textbf{(notebook, \S1--2)} on call prices. Finally,
  \textbf{show that} the argument does not extend to a payoff of the form $g(S_{T_1}, S_{T_2})$: which object, not
  constrained by (R4), then comes in? \emph{(Hint: see (R6).)}

  \item \team\ \textbf{(Notebook, \S3)} Price under M1, M2 and M3 the \emph{European-barrier put} alone (strike 100\,\%,
  barrier $B_p$ of your term sheet, maturity 3 years). Comment on the gap between M2 and M3 on
  the one hand, and on the gap with M1 on the other. What do you conclude about the role of
  the model choice for a payoff whose value depends only on $S_T$?

  \item \team\ \textbf{(Notebook, \S4)} Now price the \textbf{full Phoenix} under the three models and compute the
  break-even coupon of each. Express the M2/M3 gap in basis points per year.

  \item \team\ Where does this gap come from, given that it was essentially zero in Q11 (the put alone)?
  Identify precisely the feature of the payoff responsible, then the \textbf{two channels} through which it
  acts: (a) what the product becomes after a date on which it is not called; (b) the probability of
  being called on each date, hence the expected life of the product (compare the annuities printed
  by the notebook).

  \item \team\ \textbf{(Notebook, \S4)} Redo the computation with a protection barrier \emph{observed at each close} (see
  Section 1.1), and no longer at maturity only. By how much does the break-even coupon change?
  Does the M2--M3 gap widen? What calibration does this product require: exact at maturity
  only, or smooth over the whole surface (strikes and maturities)? \emph{(Hint: how is the vega spread
  over time?)} Finally, what must the term sheet specify?

  \item \indiv\ A structurer prices the product in local volatility and quotes the corresponding coupon. The
  structurer wins the deal against a competitor who prices in stochastic volatility. Describe the
  structurer's economic position. Is it a gain?

  \item \team\ \textbf{(Notebook, \S4bis)} We measure the \emph{forward smile} 1\,year $\to$ 3\,years in the sense of (R6),
  under M2 and under M3, together with the spot smile of the same length (from 0 to 2 years).
  The \emph{slope} is the implied vol at strike 70\,\% minus the implied vol at strike 110\,\%, in vol points.
  Compare the three slopes, and say why they are the same for every term sheet. \textbf{Deduce} the
  precise mechanism behind the gap found in Q12, and explain why calibration to vanillas does
  not detect it. \emph{(Hint: after a date on which the product is not called, what does the remaining
  product look like?)}

  \item \opt\ Model M4 (R2) is the one desks actually use. \textbf{Using Gyöngy's theorem (R4)}, show that
  exact calibration to vanillas requires
  \[
    L^2(S,t)\,\E[v_t \mid S_t = S] \;=\; \sigma^2_{\mathrm{loc}}(S,t).
  \]
  \textbf{Check that} $\eta = 0$ gives back M2 and $\eta = 1$ gives back M3. Why does this condition require
  calibrating $L$ and the simulation \emph{simultaneously} (the so-called particle method)?

  \item \team\ \textbf{(Notebook, \S4ter) Validate M4 before using it.} M4 is provided calibrated. For $\eta \in
  \{0;\, 0.25;\, 0.5;\, 0.75;\, 1\}$: (a) check that the vanillas are reproduced in every case; (b) \textbf{check the
  bounds}: does $\eta = 0$ give back the coupon of M2, and $\eta = 1$ that of M3? Quantify the gap at
  each bound and say whether it comes from the model or from the code; (c) record the break-even
  coupon and the forward smile slope, then give the \textbf{coupon range} as $\eta$ spans $[0,1]$, with
  its uncertainty. \textbf{What do you conclude} about whether $\eta$ can be identified from market data?
  \emph{The operational definition of the model uncertainty adjustment that follows, and why it is more
  satisfactory than the M2/M3 gap of Q12, go in the} Decision \emph{tab (Q36).}
\end{questions}

%======================================================================
\section{The sign of vega and skew: the trap}

\begin{questions}[resume]
  \item \indiv\ The seller of the Phoenix is short the digital strip and long the barrier put. Reasoning on
  where the forward stands relative to the coupon barrier, determine the sign of the volatility
  sensitivity of \emph{each} of the two components.

  \item \indiv\ \textbf{(Notebook, \S5)} Compute numerically the vega of each building block, then of the whole
  product. Does the result confirm the intuition? What happens to the sign if the coupon barrier
  is raised from 70\,\% to 95\,\%?

  \item \indiv\ Same question for the skew. What \emph{pricing} practice does a prudent trader adopt when the net
  sign of a parameter is not stable?
\end{questions}

%======================================================================
\section{Choosing a model by the faithfulness of its sensitivities}

So far, models have been judged on their \emph{prices}. We now turn the question around: a model is not
a price generator, it is a generator of \textbf{hedging strategies}. The right model is the one whose greeks
allow the payoff to be replicated.

\begin{questions}[resume]
  \item \team\ \textbf{(Notebook, \S6)} Assume that ``reality'' follows the reference model M3. Delta-hedge the
  European-barrier put, rebalancing every week for 3 years, with the deltas of M1, M2 and M3 in
  turn. Compare the mean, the standard deviation and the 5\,\% quantile of the terminal hedging
  P\&L.

  \item \team\ Interpret the \emph{mean} and the \emph{standard deviation} of the P\&L separately. Which one measures a
  price error, which one a greeks error? \emph{(Discussed in class, orally.)}

  \item \team\ M2 and M3 show the same initial price. Does their hedging performance tell them apart? What
  dominates the residual risk? Two leads: (a) the M3 delta in the notebook is computed at fixed
  variance $v$; when $\rho \neq 0$, is it the delta that minimizes the variance of the P\&L (see the market
  regime and the SSR, Session 4)? (b) would this residual risk disappear with a better model,
  with other hedging instruments (vanillas), or with neither? Deduce whether it is a ``model'' risk
  or whether it comes from market incompleteness when hedging with the underlying alone.

  \item \team\ \textbf{(Notebook, \S7, run by your team)} Plot the delta of the Phoenix around the protection
  barrier $B_p$ at $t = 2.5;\ 2.9;\ 2.99;\ 2.999$ (unit: see (R8)). Comment. Can a desk execute this
  hedge? This is the \emph{gap risk} of (R8). \emph{Validator's remark:} the value of the delta very close to
  maturity depends on the computational grid; why should you distrust it, and which value would
  you retain?

  \item \team\ \textbf{(Notebook, \S7)} Replace each barrier indicator with a call spread of half-width $\varepsilon$ (\emph{barrier
  shift}). Record the maximum delta as a function of $\varepsilon$ (notebook table). Then compute the effect
  of $\varepsilon$ on the break-even coupon. What do you observe, and why does \emph{symmetric} smoothing
  generate almost no margin? How does a desk actually proceed? \emph{(Hint: on which side should
  the barrier be shifted to be prudent?)}

  \item \opt\ Also plot the delta around the autocall barrier just before an observation date. What happens
  and how should it be interpreted economically?
\end{questions}

\begin{navybox}{Caveat to state explicitly}
Faithful greeks are a \textbf{necessary but not sufficient} criterion, for at least three reasons: (a) a
``reality'' has to be assumed to assess them, and that is exactly the unknown; (b) part of the
residual risk cannot be reduced by a better model: the volatility left unhedged when hedging
with the underlying alone (Q24), and the gap risk near the barrier (Q25); (c) a model can hedge
well and price badly, or the reverse. The hedging criterion therefore does not replace the pricing
criterion: it complements it.
\end{navybox}

%======================================================================
\section{Front office, independent control, prudent value}

The same product is valued at several levels, by functions with different goals and under different
frameworks. The point not to miss: the \textbf{accounting} and the \textbf{prudential} layers do not correct the
same thing and are not booked in the same place (R7).

\begin{questions}[resume]
  \item \team\ Complete and comment on the following table.
  \begin{center}
  \small
  \begin{tabular}{@{}l c c c@{}}
  \toprule
  & \textbf{Front office} & \textbf{Control (IPV)} & \textbf{Prudent value} \\
  \midrule
  Main objective & ? & ? & ? \\
  Framework & ? & ? & ? \\
  Success criterion & ? & ? & ? \\
  Frequency & ? & ? & ? \\
  Treatment of barriers & ? & ? & ? \\
  Attitude toward complexity & ? & ? & ? \\
  \bottomrule
  \end{tabular}
  \end{center}

  \item \team\ \textbf{Using (R7)}, explain why IPV is \emph{not} a reserve, and what it produces instead. What happens
  when an IPV difference is found and judged significant?

  \item \team\ Distinguish precisely a \textbf{fair value adjustment} from an \textbf{AVA}: (a) under which framework
  is each one required; (b) what does each one try to approximate --- exit price? prudent value,
  and at what confidence level?; (c) where is each one booked; (d) AVAs are added on top of
  the accounting adjustments already booked: how does one make sure that the same risk is not
  counted twice?

  \item \team\ \textbf{Classify} each of the uncertainties met in this tutorial according to the three kinds of (R7) ---
  poorly observed market parameter, model parameter, specification --- then state the adjustment
  that most naturally deals with it.
  \begin{center}
  \small
  \begin{tabularx}{\linewidth}{@{}>{\raggedright\arraybackslash}X c c c@{}}
  \toprule
  Uncertainty & Kind & Accounting adjustment & AVA category \\
  \midrule
  Mixing fraction $\eta$ (Q18) & ? & ? & ? \\
  Choice local vol / local-stochastic & ? & ? & ? \\
  Projected dividends at 3 years (single stock) & ? & ? & ? \\
  Implied correlation of a \emph{custom} basket & ? & ? & ? \\
  Volatility left unhedged by the delta alone (Q24) & ? & ? & ? \\
  Gap risk at the barrier (Q25) & ? & ? & ? \\
  Gap between theoretical and smoothed price (Q26) & ? & ? & ? \\
  \bottomrule
  \end{tabularx}
  \end{center}

  \item \team\ The front office smooths its barriers (\emph{barrier shift}) and hedges the regularized product. Should
  independent control validate this value? Argue both ways, then say which layer deals with the
  remaining gap.

  \item \team\ The front office prices in local volatility; the independent valuation team uses M4. The gap is
  persistent. \textbf{How does it show}, first in the accounting result and then in capital? Why is it
  not the same amount?

  \item \team\ A product is sold at 100 while the validated model values it at 98. What happens to the gap
  of 2 in the accounts if the product is classified as Level 3 in the fair value hierarchy? How does
  this rule protect against the situation described in Q15?

  \item \team\ State the central tension in one sentence, contrasting \emph{accuracy of the representation} and \emph{operability of the hedge}. Can it be settled?

  \item \team\ A desk proposes to replace its local volatility model with model M4, for which Q18 measured
  the coupon range as the mixing fraction $\eta$ spans $[0,1]$ and assessed whether $\eta$ can be identified
  from vanillas. Should the validation committee approve? Which safeguards should it require,
  on what principle should $\eta$ be set, and to which layer of (R7) should this range be attached?
\end{questions}

%======================================================================
\section{Counter-example: when \emph{not} to enrich the model}

\begin{greybox}{Reminder: what the product is, economically}
Before tackling this part, keep in mind that the Phoenix is not only a payoff. It is a \textbf{debt security}
issued on the issuer's funding curve, \textbf{redeemed at a random date} that depends on the path of
the underlying, sold at a fixed coupon during several weeks of marketing, and then quoted by a
single market maker. Its life is not a parameter: it is a random variable. This observation alone
opens a dimension of risk that everything above has ignored.
\end{greybox}

\begin{questions}[resume]
  \item \indiv\ The seller hedges the rate leg with zero-coupon bonds. To isolate the mechanism, take the
  simplest autocall: \textbf{maturity 2 years, with a single early-redemption date at 1 year} (the
  note is called at 1 year if $S_1 \geq B_a S_0$, and otherwise runs to its maturity at 2 years, where it is
  redeemed). Until the 1-year date, the issuer does not know whether it will repay after 1 year or
  after 2 years: its rate hedge is a mix of 1-year and 2-year zero-coupon bonds, weighted by the
  probability of being called. Describe how this rate hedge must be rebalanced in the following
  two cases: (a) the underlying rises; (b) the underlying falls. Then say what happens to the
  hedge \emph{at} the 1-year date if the note is called, and \emph{between 1 and 2 years} if it is not.

  \item \indiv\ Assuming a \emph{positive} correlation between the underlying and interest rates, show that the
  rebalancing loses money in both cases. Redo the reasoning with a negative correlation. Deduce
  the direction of the effect on the price.

  \item \indiv\ How does this cost vary with the maturity of the product? With the number of early-redemption dates?

  \item \indiv\ Should one therefore adopt a stochastic interest rate model? Build the argument \emph{against},
  carefully distinguishing what a model lets you \emph{see} from what it lets you \emph{do}.

  \item \opt\ What, then, is market practice? \emph{(Practitioner's order of magnitude, given for guidance only:
  about 10 bp of coupon per year, depending on the maturity and the size of the deal.)}

  \item \team\ \textbf{Synthesis.} State the general criterion that emerges from the whole tutorial: under what
  condition is enriching a model justified? Apply this criterion to the four decisions met along
  the way: (a) moving from Black--Scholes to local volatility; (b) moving from local volatility to
  stochastic volatility; (c) moving to local-stochastic M4; (d) moving to stochastic rates. \emph{In the
  Decision tab, only (a) to (c) are graded; (d) draws on the individual questions Q37--Q40.}
\end{questions}

%======================================================================
\section{Block D --- A pricer written by an AI: validate it}

\begin{greybox}{Why this block}
A code assistant produces in a few seconds a clean, documented pricer that runs and prints a
plausible coupon. The cost of \emph{production} has collapsed; the cost of \emph{verification} has not moved.
The job shifts toward designing tests: this is exactly form (i) of the taxonomy of Section 3, which
the rest of the tutorial sets aside.
\end{greybox}

\noindent\textbf{Organization.} At the start of TD 1, each team designates a \textbf{Block D lead} and enters his or her
email address in the answer sheet (tab \emph{Team}). Once the team has handed in its test plan (D1), the
instructor emails the team's module, \texttt{phoenix\_pricer\_llm.py}, to the Block D lead, who shares it
with the rest of the team. The module was generated by an LLM from the description of \emph{your} term
sheet and of the reference market. It contains several errors. You are the validator. Modules differ
from one team to another and from one tutorial group to another.

\begin{greybox}{Interface of the module (all you need to write the plan)}
\small
\texttt{simulate\_heston(S0, T, n\_steps, n\_paths, r, q, hp, seed=None)}: spot paths, an array of shape
$(n_{\text{steps}}+1) \times n_{\text{paths}}$ (one column per path).\\
\texttt{phoenix\_pv(paths, T, coupon, product, r)}: discounted value of the cash flows, path by path (notional
1); \texttt{paths} has the same shape as above, and hand-built paths can be passed in.\\
\texttt{price(coupon, S0=None, n\_steps=300, n\_paths=50\_000, seed=None, \ldots)} (pass arguments by name:
the second positional argument is the spot), \texttt{break\_even\_coupon(\ldots)}, \texttt{delta(coupon, bump=0.001, \ldots)}.
Constants \texttt{PRODUCT}, \texttt{MARKET}, \texttt{HESTON}. The coupon is \emph{semi-annual}. The module calls the vol of vol \texttt{xi} (the
notebook calls it \texttt{sigma}): do not pass the notebook's parameters to the module directly.
\end{greybox}

\begin{dquestions}
  \item \team\ \textbf{Before receiving the module}, write a test plan (half a page), handed in on Moodle or on
  paper: \textbf{at least five tests, one per error class} (numerical scheme, market data, payoff logic,
  term-sheet reading, greeks), each with its expected result, justified without the module. Hint:
  a calibrated model must reproduce a martingale, known vanilla prices, and hand-built paths; a
  greek must be stable from one computation to the next.

  \item \team\ Run your tests against the module. List those that fail, and locate the line responsible. One
  error can hide another: once you have found one, fix it and rerun your tests.

  \item \team\ For each error found: fix it, then measure its impact in basis points of annual break-even
  coupon, \emph{all else being equal} (same random draws).

  \item \team\ The assistant claims that ``the break-even coupon it returns is consistent with the literature.''
  What has it actually checked? What should have been checked?

  \item \team\ Is this pricer a \emph{model} in the sense of definition (R9), even though it was written by a generative
  AI tool? Who must validate it, under which change category (I to III, see (R9)), and how can
  one prevent such a tool from circulating outside the model inventory (\emph{shadow IT})? \emph{(Everything
  you need is in (R9); no external text is required.)}
\end{dquestions}

\begin{navybox}{Rule of the game}
Block D is the only part where an AI may be used, and only to help write your tests. You remain
solely responsible for what you hand in: the \emph{Decision} tab states which tools were used for the tests
and what you checked yourselves. Everywhere else, AI tools are not allowed.
\end{navybox}

%======================================================================
\section{The \emph{Decision} tab (TD 2, team)}

There is no oral presentation: the team's validation decision is written in the \emph{Decision} tab of the team
answer sheet. It contains, in short answers:
\begin{itemize}
  \item \textbf{the decision}, for the model \emph{and} for the Block D pricer: approve, approve with conditions, or
  reject, justified in five lines with the numbers of your term sheet and their precision, and what
  would change it;
  \item \textbf{the model uncertainty adjustment}: the $\eta$ range retained (Q18) and the principle used to set
  $\eta$ (Q36);
  \item \textbf{Section 7}: Q28 to Q35;
  \item \textbf{Block D}: D4 and D5 (the impacts of D3 are in the \emph{Block D} tab);
  \item \textbf{the criterion} of Q42, parts (a) to (c);
  \item \textbf{Block D AI statement}: the AI tools used to write your Block D tests, if any, and what you
  checked yourselves (or ``none'').
\end{itemize}
Start it in TD 2 (minutes 50--75); the team has one week after TD 2 to finish and submit it.

\section*{Appendix --- The two answer sheets}

The answer sheet is an \textbf{Excel file}, \texttt{TD\_answer\_sheet.xlsx}, posted on the ENT. One row per number,
for \emph{your} term sheet; add rows if needed. It is filled in during both tutorials and submitted on Moodle
one week after TD 2, together with the notebook. Its tabs: \emph{Team} (group, team number, members,
Block D lead's email), \emph{Answers} (question, numbers, paths, standard error, one-sentence reading),
\emph{Block D} (test plan and results), \emph{Decision} (Section 10), \emph{Teams} (term sheet of every team).

The \textbf{individual answer sheet}, \texttt{TD\_individual\_answers.xlsx}, has one cell per \indiv\ question. Write
formulas in plain text (e.g.\ \texttt{V0 = E[sum exp(-r t\_i) c 1\{R\_i >= Bc\} ...]}); long derivations may
go in an optional PDF appendix, submitted with it.

\vspace{2em}
\noindent\rule{\textwidth}{0.4pt}
{\small\emph{Useful references}: M. Bouzoubaa \& A. Osseiran, \emph{Exotic Options and Hybrids}, Wiley, 2010 (ch.\ 4, 10, 11, 12);
J. Gatheral, \emph{The Volatility Surface}, Wiley, 2006; L. Bergomi, \emph{Stochastic Volatility Modeling}, CRC, 2016; Basel
Committee / CRR art.\ 105 on \emph{prudent valuation}; J. Guyon \& P. Henry-Labordère, \emph{Nonlinear Option Pricing},
CRC, 2014 (local-stochastic model, particle method).\par}

\end{document}
